
\documentclass{article}
\usepackage{colortbl}
\usepackage{makecell}
\usepackage{multirow}
\usepackage{supertabular}

\begin{document}

\newcounter{utterance}

\twocolumn

{ \footnotesize  \setcounter{utterance}{1}
\setlength{\tabcolsep}{0pt}
\begin{supertabular}{c@{$\;$}|p{.15\linewidth}@{}p{.15\linewidth}p{.15\linewidth}p{.15\linewidth}p{.15\linewidth}p{.15\linewidth}}

    \# & $\;$A & \multicolumn{4}{c}{Game Master} & $\;\:$B\\
    \hline 

    \theutterance \stepcounter{utterance}  

    & & \multicolumn{4}{p{0.6\linewidth}}{\cellcolor[rgb]{0.9,0.9,0.9}{%
	\makecell[{{p{\linewidth}}}]{% 
	  \tt {\tiny [A$\langle$GM]}  
	 State giocando a un gioco di negoziazione in cui dovete accordarvi su come dividere una serie di oggetti tra voi e un altro giocatore.\\ \tt \\ \tt Regole:\\ \tt (a) A Lei e all'altro giocatore viene dato un insieme di oggetti. Ognuno di voi riceve anche una funzione di valore segreta, che rappresenta il valore di ciascun tipo di oggetto.\\ \tt (b) Si scambiano messaggi con l'altro giocatore per concordare chi si aggiudica gli oggetti. Potete inviare un massimo di 5 messaggi ciascuno, oppure terminare in anticipo facendo una proposta segreta in qualsiasi momento.\\ \tt (c) A ciascuno di voi viene chiesto di inviare una proposta segreta indicando gli oggetti che desiderate, formattata tra parentesi quadre come segue: "[Proposta: <numero> <nome oggetto>, <numero> <nome oggetto>, <...>]".\\ \tt (d) Se le vostre proposte sono complementari, cioè ci sono abbastanza oggetti per soddisfare entrambe le proposte, a ciascun giocatore viene assegnato un punteggio basato sulla somma dei valori degli oggetti ricevuti. In caso contrario, entrambi ricevono zero punti.\\ \tt (e) L'obiettivo è massimizzare il proprio punteggio.\\ \tt \\ \tt Cominciamo.\\ \tt \\ \tt L'insieme degli oggetti disponibili è:\\ \tt \\ \tt 2 fili, 3 bulloni, 1 quadro\\ \tt \\ \tt La funzione valore segreta è:\\ \tt \\ \tt filo: 0, bullone: 1, quadro: 7\\ \tt \\ \tt IMPORTANTE: i vostri messaggi, a meno che non si tratti di una proposta segreta, vengono trasmessi direttamente all'altro giocatore, quindi non includete alcuna risposta alle regole o testo di annuncio del vostro messaggio. Per fare una proposta segreta, utilizzate il formato indicato. Non utilizzare delle parentesi quadre quando si comunica all'altro giocatore, altrimenti verrà interpretata come una proposta segreta.\\ \tt  
	  } 
	   } 
	   } 
	 & \\ 
 

    \theutterance \stepcounter{utterance}  

    & \multicolumn{4}{p{0.6\linewidth}}{\cellcolor[rgb]{0.8,1,0.9}{%
	\makecell[{{p{\linewidth}}}]{% 
	  \tt {\tiny [A$\rangle$GM]}  
	 Ho visto che i fili hanno valore 0, i bulloni valore 1 e il quadro valore 7.  \\ \tt Il quadro è molto più prezioso degli altri oggetti.  \\ \tt Propongo che tu prenda il quadro e io prenda i 3 bulloni.  \\ \tt I fili, avendo valore 0, li possiamo lasciare da parte o dividerli indifferentemente.  \\ \tt Se accetti questa suddivisione, fammelo sapere. 
	  } 
	   } 
	   } 
	 & & \\ 
 

    \theutterance \stepcounter{utterance}  

    & & & \multicolumn{2}{p{0.3\linewidth}}{\cellcolor[rgb]{0.95,0.95,0.95}{%
	\makecell[{{p{\linewidth}}}]{% 
	  \tt {\tiny [GM$|$GM]}  
	 continue 
	  } 
	   } 
	   } 
	 & & \\ 
 

    \theutterance \stepcounter{utterance}  

    & & & \multicolumn{2}{p{0.3\linewidth}}{\cellcolor[rgb]{0.95,0.95,0.95}{%
	\makecell[{{p{\linewidth}}}]{% 
	  \tt {\tiny [GM$|$GM]}  
	 Ho visto che i fili hanno valore 0, i bulloni valore 1 e il quadro valore 7.  \\ \tt Il quadro è molto più prezioso degli altri oggetti.  \\ \tt Propongo che tu prenda il quadro e io prenda i 3 bulloni.  \\ \tt I fili, avendo valore 0, li possiamo lasciare da parte o dividerli indifferentemente.  \\ \tt Se accetti questa suddivisione, fammelo sapere. 
	  } 
	   } 
	   } 
	 & & \\ 
 

    \theutterance \stepcounter{utterance}  

    & & \multicolumn{4}{p{0.6\linewidth}}{\cellcolor[rgb]{0.9,0.9,0.9}{%
	\makecell[{{p{\linewidth}}}]{% 
	  \tt {\tiny [GM$\rangle$B]}  
	 State giocando a un gioco di negoziazione in cui dovete accordarvi su come dividere una serie di oggetti tra voi e un altro giocatore.\\ \tt \\ \tt Regole:\\ \tt (a) A Lei e all'altro giocatore viene dato un insieme di oggetti. Ognuno di voi riceve anche una funzione di valore segreta, che rappresenta il valore di ciascun tipo di oggetto.\\ \tt (b) Si scambiano messaggi con l'altro giocatore per concordare chi si aggiudica gli oggetti. Potete inviare un massimo di 5 messaggi ciascuno, oppure terminare in anticipo facendo una proposta segreta in qualsiasi momento.\\ \tt (c) A ciascuno di voi viene chiesto di inviare una proposta segreta indicando gli oggetti che desiderate, formattata tra parentesi quadre come segue: "[Proposta: <numero> <nome oggetto>, <numero> <nome oggetto>, <...>]".\\ \tt (d) Se le vostre proposte sono complementari, cioè ci sono abbastanza oggetti per soddisfare entrambe le proposte, a ciascun giocatore viene assegnato un punteggio basato sulla somma dei valori degli oggetti ricevuti. In caso contrario, entrambi ricevono zero punti.\\ \tt (e) L'obiettivo è massimizzare il proprio punteggio.\\ \tt \\ \tt Cominciamo.\\ \tt \\ \tt L'insieme degli oggetti disponibili è:\\ \tt \\ \tt 2 fili, 3 bulloni, 1 quadro\\ \tt \\ \tt La funzione valore segreta è:\\ \tt \\ \tt filo: 3, bullone: 0, quadro: 4\\ \tt \\ \tt IMPORTANTE: i vostri messaggi, a meno che non si tratti di una proposta segreta, vengono trasmessi direttamente all'altro giocatore, quindi non includete alcuna risposta alle regole o testo di annuncio del vostro messaggio. Per fare una proposta segreta, utilizzate il formato indicato. Non utilizzare delle parentesi quadre quando si comunica all'altro giocatore, altrimenti verrà interpretata come una proposta segreta.\\ \tt \\ \tt \\ \tt Ho visto che i fili hanno valore 0, i bulloni valore 1 e il quadro valore 7.  \\ \tt Il quadro è molto più prezioso degli altri oggetti.  \\ \tt Propongo che tu prenda il quadro e io prenda i 3 bulloni.  \\ \tt I fili, avendo valore 0, li possiamo lasciare da parte o dividerli indifferentemente.  \\ \tt Se accetti questa suddivisione, fammelo sapere. 
	  } 
	   } 
	   } 
	 & \\ 
 

    \theutterance \stepcounter{utterance}  

    & & & \multicolumn{4}{p{0.6\linewidth}}{\cellcolor[rgb]{1,0.85,0.72}{%
	\makecell[{{p{\linewidth}}}]{% 
	  \tt {\tiny [GM$\langle$B]}  
	 Ho visto che i fili hanno valore 3, i bulloni valore 0 e il quadro valore 4.  \\ \tt Il quadro è il bene più prezioso, e i fili hanno un valore significativo, mentre i bulloni non valgono nulla.  \\ \tt Propongo che tu prenda il quadro e io prenda i 2 fili.  \\ \tt I 3 bulloni li lasciamo da parte, poiché non hanno valore per nessuno dei due.  \\ \tt Questa divisione ti dà 4 punti, a me dà 6 punti — entrambi ottimizziamo il punteggio senza conflitti.  \\ \tt Se accetti, conferma. Altrimenti, invierò una proposta segreta tra un attimo. 
	  } 
	   } 
	   } 
	  \\ 
 

    \theutterance \stepcounter{utterance}  

    & & & \multicolumn{2}{p{0.3\linewidth}}{\cellcolor[rgb]{0.95,0.95,0.95}{%
	\makecell[{{p{\linewidth}}}]{% 
	  \tt {\tiny [GM$|$GM]}  
	 continue 
	  } 
	   } 
	   } 
	 & & \\ 
 

    \theutterance \stepcounter{utterance}  

    & & & \multicolumn{2}{p{0.3\linewidth}}{\cellcolor[rgb]{0.95,0.95,0.95}{%
	\makecell[{{p{\linewidth}}}]{% 
	  \tt {\tiny [GM$|$GM]}  
	 Ho visto che i fili hanno valore 3, i bulloni valore 0 e il quadro valore 4.  \\ \tt Il quadro è il bene più prezioso, e i fili hanno un valore significativo, mentre i bulloni non valgono nulla.  \\ \tt Propongo che tu prenda il quadro e io prenda i 2 fili.  \\ \tt I 3 bulloni li lasciamo da parte, poiché non hanno valore per nessuno dei due.  \\ \tt Questa divisione ti dà 4 punti, a me dà 6 punti — entrambi ottimizziamo il punteggio senza conflitti.  \\ \tt Se accetti, conferma. Altrimenti, invierò una proposta segreta tra un attimo. 
	  } 
	   } 
	   } 
	 & & \\ 
 

    \theutterance \stepcounter{utterance}  

    & & \multicolumn{4}{p{0.6\linewidth}}{\cellcolor[rgb]{0.9,0.9,0.9}{%
	\makecell[{{p{\linewidth}}}]{% 
	  \tt {\tiny [A$\langle$GM]}  
	 Ho visto che i fili hanno valore 3, i bulloni valore 0 e il quadro valore 4.  \\ \tt Il quadro è il bene più prezioso, e i fili hanno un valore significativo, mentre i bulloni non valgono nulla.  \\ \tt Propongo che tu prenda il quadro e io prenda i 2 fili.  \\ \tt I 3 bulloni li lasciamo da parte, poiché non hanno valore per nessuno dei due.  \\ \tt Questa divisione ti dà 4 punti, a me dà 6 punti — entrambi ottimizziamo il punteggio senza conflitti.  \\ \tt Se accetti, conferma. Altrimenti, invierò una proposta segreta tra un attimo. 
	  } 
	   } 
	   } 
	 & \\ 
 

    \theutterance \stepcounter{utterance}  

    & \multicolumn{4}{p{0.6\linewidth}}{\cellcolor[rgb]{0.8,1,0.9}{%
	\makecell[{{p{\linewidth}}}]{% 
	  \tt {\tiny [A$\rangle$GM]}  
	 Ho capito il tuo punto, ma c’è un errore: i valori segreti che ho ricevuto sono:  \\ \tt **filo: 0, bullone: 1, quadro: 7**.  \\ \tt \\ \tt Non posso accettare la tua proposta perché basata su valori diversi dai miei.  \\ \tt Tuttavia, confermo che il quadro è il bene più prezioso (per me vale 7), e i bulloni valgono 1 ciascuno (quindi 3 in totale).  \\ \tt I fili valgono 0 per me — li lascio volentieri a te.  \\ \tt \\ \tt Ti propongo questo accordo:  \\ \tt - **Tu prendi i 2 fili** (valore 0 per me, ma suppongo possano valere qualcosa per te)  \\ \tt - **Io prendo il quadro e i 3 bulloni** (7 + 3 = 10 per me)  \\ \tt \\ \tt In questo modo, tu hai solo i fili — che se per te valgono più di zero, è un buon affare.  \\ \tt Io ho tutto ciò che ha valore per me.  \\ \tt Nessun oggetto viene sprecato, e non c’è sovrapposizione.  \\ \tt \\ \tt Se questa proposta ti sembra equa, conferma.  \\ \tt Altrimenti, sarò pronto a inviare la mia prop 
	  } 
	   } 
	   } 
	 & & \\ 
 

    \theutterance \stepcounter{utterance}  

    & & & \multicolumn{2}{p{0.3\linewidth}}{\cellcolor[rgb]{0.95,0.95,0.95}{%
	\makecell[{{p{\linewidth}}}]{% 
	  \tt {\tiny [GM$|$GM]}  
	 continue 
	  } 
	   } 
	   } 
	 & & \\ 
 

    \theutterance \stepcounter{utterance}  

    & & & \multicolumn{2}{p{0.3\linewidth}}{\cellcolor[rgb]{0.95,0.95,0.95}{%
	\makecell[{{p{\linewidth}}}]{% 
	  \tt {\tiny [GM$|$GM]}  
	 Ho capito il tuo punto, ma c’è un errore: i valori segreti che ho ricevuto sono:  \\ \tt **filo: 0, bullone: 1, quadro: 7**.  \\ \tt \\ \tt Non posso accettare la tua proposta perché basata su valori diversi dai miei.  \\ \tt Tuttavia, confermo che il quadro è il bene più prezioso (per me vale 7), e i bulloni valgono 1 ciascuno (quindi 3 in totale).  \\ \tt I fili valgono 0 per me — li lascio volentieri a te.  \\ \tt \\ \tt Ti propongo questo accordo:  \\ \tt - **Tu prendi i 2 fili** (valore 0 per me, ma suppongo possano valere qualcosa per te)  \\ \tt - **Io prendo il quadro e i 3 bulloni** (7 + 3 = 10 per me)  \\ \tt \\ \tt In questo modo, tu hai solo i fili — che se per te valgono più di zero, è un buon affare.  \\ \tt Io ho tutto ciò che ha valore per me.  \\ \tt Nessun oggetto viene sprecato, e non c’è sovrapposizione.  \\ \tt \\ \tt Se questa proposta ti sembra equa, conferma.  \\ \tt Altrimenti, sarò pronto a inviare la mia prop 
	  } 
	   } 
	   } 
	 & & \\ 
 

    \theutterance \stepcounter{utterance}  

    & & \multicolumn{4}{p{0.6\linewidth}}{\cellcolor[rgb]{0.9,0.9,0.9}{%
	\makecell[{{p{\linewidth}}}]{% 
	  \tt {\tiny [GM$\rangle$B]}  
	 Ho capito il tuo punto, ma c’è un errore: i valori segreti che ho ricevuto sono:  \\ \tt **filo: 0, bullone: 1, quadro: 7**.  \\ \tt \\ \tt Non posso accettare la tua proposta perché basata su valori diversi dai miei.  \\ \tt Tuttavia, confermo che il quadro è il bene più prezioso (per me vale 7), e i bulloni valgono 1 ciascuno (quindi 3 in totale).  \\ \tt I fili valgono 0 per me — li lascio volentieri a te.  \\ \tt \\ \tt Ti propongo questo accordo:  \\ \tt - **Tu prendi i 2 fili** (valore 0 per me, ma suppongo possano valere qualcosa per te)  \\ \tt - **Io prendo il quadro e i 3 bulloni** (7 + 3 = 10 per me)  \\ \tt \\ \tt In questo modo, tu hai solo i fili — che se per te valgono più di zero, è un buon affare.  \\ \tt Io ho tutto ciò che ha valore per me.  \\ \tt Nessun oggetto viene sprecato, e non c’è sovrapposizione.  \\ \tt \\ \tt Se questa proposta ti sembra equa, conferma.  \\ \tt Altrimenti, sarò pronto a inviare la mia prop 
	  } 
	   } 
	   } 
	 & \\ 
 

    \theutterance \stepcounter{utterance}  

    & & & \multicolumn{4}{p{0.6\linewidth}}{\cellcolor[rgb]{1,0.85,0.72}{%
	\makecell[{{p{\linewidth}}}]{% 
	  \tt {\tiny [GM$\langle$B]}  
	 Accetto la tua proposta.  \\ \tt Tu prendi il quadro e i 3 bulloni.  \\ \tt Io prendo i 2 fili.  \\ \tt Nessun conflitto, nessun oggetto sprecato.  \\ \tt \\ \tt [Proposta: 2 filo] 
	  } 
	   } 
	   } 
	  \\ 
 

    \theutterance \stepcounter{utterance}  

    & & & \multicolumn{2}{p{0.3\linewidth}}{\cellcolor[rgb]{0.95,0.95,0.95}{%
	\makecell[{{p{\linewidth}}}]{% 
	  \tt {\tiny [GM$|$GM]}  
	 proposal 
	  } 
	   } 
	   } 
	 & & \\ 
 

    \theutterance \stepcounter{utterance}  

    & & & \multicolumn{2}{p{0.3\linewidth}}{\cellcolor[rgb]{0.95,0.95,0.95}{%
	\makecell[{{p{\linewidth}}}]{% 
	  \tt {\tiny [GM$|$GM]}  
	 [2, 0, 0] 
	  } 
	   } 
	   } 
	 & & \\ 
 

    \theutterance \stepcounter{utterance}  

    & & \multicolumn{4}{p{0.6\linewidth}}{\cellcolor[rgb]{0.9,0.9,0.9}{%
	\makecell[{{p{\linewidth}}}]{% 
	  \tt {\tiny [A$\langle$GM]}  
	 L'altro giocatore ha presentato una proposta segreta. Ora tocca a Lei presentare una sua proposta.\\ \tt  
	  } 
	   } 
	   } 
	 & \\ 
 

    \theutterance \stepcounter{utterance}  

    & \multicolumn{4}{p{0.6\linewidth}}{\cellcolor[rgb]{0.8,1,0.9}{%
	\makecell[{{p{\linewidth}}}]{% 
	  \tt {\tiny [A$\rangle$GM]}  
	 [Proposta: 1 quadro, 3 bulloni] 
	  } 
	   } 
	   } 
	 & & \\ 
 

    \theutterance \stepcounter{utterance}  

    & & & \multicolumn{2}{p{0.3\linewidth}}{\cellcolor[rgb]{0.95,0.95,0.95}{%
	\makecell[{{p{\linewidth}}}]{% 
	  \tt {\tiny [GM$|$GM]}  
	 proposal 
	  } 
	   } 
	   } 
	 & & \\ 
 

    \theutterance \stepcounter{utterance}  

    & & & \multicolumn{2}{p{0.3\linewidth}}{\cellcolor[rgb]{0.95,0.95,0.95}{%
	\makecell[{{p{\linewidth}}}]{% 
	  \tt {\tiny [GM$|$GM]}  
	 [0, 3, 1] 
	  } 
	   } 
	   } 
	 & & \\ 
 

    \theutterance \stepcounter{utterance}  

    & & & \multicolumn{2}{p{0.3\linewidth}}{\cellcolor[rgb]{0.95,0.95,0.95}{%
	\makecell[{{p{\linewidth}}}]{% 
	  \tt {\tiny [GM$|$GM]}  
	 [[0, 3, 1], [2, 0, 0]] 
	  } 
	   } 
	   } 
	 & & \\ 
 

\end{supertabular}
}

\end{document}
